\clearpage
\section{A combined contextual copying subsystem}
\label{sec:copying-combined}
The previously separate allocation, dereference, assignment, atom-equality,
application and copying machines now execute within one transition system.
Every proved local run embeds with precisely its original cost. Graph-changing
transitions can occur in shallow typed contexts, including an argument of an
active operation. A concrete allowance-indexed execution instance transports
all finite runs under recursive graph isomorphism. Its execution-isomorphism
law, previously an abstract premise, is proved; the existing graph-context
congruence theorem therefore applies to this instance without that premise.

This is a subsystem milestone on the route to Proposition~12.2. The original
terms, Figure~16 translation, stable-argument premises and copyability premises
are unchanged. Section~\ref{sec:copying-arithmetic-padding} adds arithmetic and
forward exact-cost padding to this same machine. Matching independence, deterministic
behaviour up to representation, refocusing, focused-graph correspondence,
reverse padding transport and specimen locality are not proved for this combined
system. The named representation relation admits isomorphic presentations of
each rewrite; this is an explicit semantic closure, not a theorem asserting
unique raw graph matches. The final proposition and all template-robustness
obligations remain open.

\Needspace{12\baselineskip}
\subsection{Local rules and shallow matching}

\termdef{combined-copying-subsystem}
A local rule records two actual finite hypernets with the same ordered
interface and one active coordinate on each side. Its constructors are reference creation, dereference, assignment, same-atom
equality, different-atom equality, micro-beta and copying, together with the
unary and binary arithmetic constructors added in
Section~\ref{sec:copying-arithmetic-padding}. Creation, assignment and
micro-beta retain stability of the supplied value; copying retains copyability
and its arbitrary residual contraction tree. The two equality constructors
retain their separate finite graphs. The coordinates are their original term
inputs. No constructor assumes termination or the desired equivalence.

\lean{Copying/CombinedRules}{10}{36}{Combined.Local}

\termdef{shallow-rewrite-match}
The rewrite relation is the least relation generated from these local pairs
by three further constructors. Recursive isomorphisms may rename both sides
and their active vertices independently. Typed replacement may put both sides
in the same selected edge slot of a host graph. Simultaneous expansion may
change one selected component while every other filling remains equal. The
selected host edge is erased by replacement: its label supplies the slot's
signature and is not asserted to be an observable hole label. These
constructors give explicit witnesses for a match; they do not implement a
unique-match algorithm.

All insertions take place at the outer graph level. There is no constructor
that evaluates inside an unopened box. A recursive isomorphism may rename
vertices inside a box, but it preserves the box hierarchy and does not move
the token into its body. Box opening is the actual micro-beta rule. Later
binding-free template arguments must still justify their own restrictions;
shallow closure alone does not prove template robustness.

\lean{Copying/CombinedRules}{41}{59}{Combined.Rewrite}

At fixed allowance $d$ and interfaces $I,O$, a state is a dependent pair
of an actual hypernet and a search position on that graph. A transition is
either the existing intrinsic search relation on one unchanged graph, or a
rewrite from rewrite focus at the source coordinate to search focus at the
target coordinate. Finality is backtracking at an external input. Search
cannot continue there, and a rewrite requires a different focus mode;
case analysis therefore proves final normality and supplies a machine.

\lean{Copying/CombinedRules}{61}{61}{Combined.State}

\lean{Copying/CombinedRules}{64}{66}{Combined.Step}

\lean{Copying/CombinedRules}{68}{76}{Combined.final_normal}

\lean{Copying/CombinedRules}{78}{81}{Combined.machine}

Every local rule becomes a transition by the local-match constructor
followed by the rewrite-step constructor. Conversely, rewrite focus has no
intrinsic search successor, by the previously checked search theorem. Thus
copying and compute are explicit graph-changing transitions; neither has
been put into the older linear search or compute parameter.

\lean{Copying/CombinedRules}{84}{85}{Combined.local_step}

\lean{Copying/CombinedRules}{88}{89}{Combined.rewrite_not_search}

\Needspace{12\baselineskip}
\subsection{State isomorphisms and exact costs}

\termdef{combined-state-isomorphism}
A state isomorphism consists of a recursive graph isomorphism, equality of
focus modes and equality of the target token coordinate with the image of
the source coordinate. Reflexivity uses identity bijections. Symmetry uses
the inverse vertex map. Transitivity composes graph isomorphisms and vertex
maps; the port-map composition equation follows by the three cases of a
port. These maps fix ordered external inputs and outputs. Consequently a
token is at an input before transport exactly when it is at an input after
transport, proving equivalence of finality.

\lean{Copying/CombinedIso}{10}{13}{Combined.StateIso}

\lean{Copying/CombinedIso}{17}{20}{Combined.StateIso.refl}

\lean{Copying/CombinedIso}{21}{26}{Combined.StateIso.symm}

\lean{Copying/CombinedIso}{27}{35}{Combined.StateIso.trans}

\lean{Copying/CombinedIso}{37}{39}{Combined.StateIso.final_iff}

One-step transport is proved by cases on the transition. Search uses
the established incidence transport through the induced boundary isomorphism.
For a rewrite, the explicit renaming closure changes its source presentation;
the old target itself can be retained, with the reflexive target isomorphism.
In both cases there is one target transition and an isomorphism relating
endpoints. This proof uses the chosen rewrite relation's representation
closure; it does not prove that arbitrary redex decompositions are unique.

\lean{Copying/CombinedIso}{43}{65}{Combined.step_transport}

Induction on a run transports one step and then the remainder using
the newly obtained endpoint isomorphism. No stuttering steps are inserted,
so the length stays exactly $n$. Successful termination at cost $n$ means
existence of a final endpoint after $n$ transitions. Run transport and finality
transport prove the forward implication; the inverse state isomorphism proves
the converse. No determinism assumption is used.

\lean{Copying/CombinedIso}{68}{76}{Combined.run_transport}

\lean{Copying/CombinedIso}{79}{80}{Combined.TerminatesIn}

\lean{Copying/CombinedIso}{82}{90}{Combined.terminatesIn_iso}

\Needspace{12\baselineskip}
\subsection{Changing graphs inside an active context}

Plugging a state inserts both its graph and its token into the same
typed host slot. The two mode components agree by construction. An intrinsic
search step lifts by the established insertion theorem; a rewrite lifts by
the corresponding rewrite constructor. Induction extends this to all finite
runs with the exact same length. This theorem concerns transitions, not
preservation of finality at an arbitrarily attached boundary.

\lean{Copying/CombinedLocality}{10}{13}{Combined.plugState}

\lean{Copying/CombinedLocality}{15}{21}{Combined.step_plug}

\lean{Copying/CombinedLocality}{23}{30}{Combined.run_plug}

Serial state insertion separately handles a changing left component
followed by a fixed continuation, and a fixed left component followed by a
changing right component. Each construction maps the token with the actual
serial insertion map. A left external input remains an external input; the
proof computes the source port of the left frame hole. One-step lifting uses
search transport or simultaneous rewrite expansion with the other component
unchanged. The unchanged-component proof is a case split on the two hole
tags. Run lifting then follows by induction.

\lean{Copying/CombinedLocality}{33}{34}{Combined.serialLeft}

\lean{Copying/CombinedLocality}{35}{36}{Combined.serialRight}

\lean{Copying/CombinedLocality}{38}{43}{Combined.left_input}

\lean{Copying/CombinedLocality}{45}{57}{Combined.step_serialLeft}

\lean{Copying/CombinedLocality}{59}{71}{Combined.step_serialRight}

\lean{Copying/CombinedLocality}{73}{78}{Combined.run_serialLeft}

\lean{Copying/CombinedLocality}{80}{85}{Combined.run_serialRight}

A right-component input becomes an internal connection of the serial
frame. Backtracking there is therefore not final for the whole graph.
The proof computes the right hole's source port and contradicts its alleged
external-input status. A left-component input instead remains external,
so backtracking there is final. These two results prevent local completion
from being incorrectly promoted to program termination.

\lean{Copying/CombinedLocality}{88}{95}{Combined.right_return_not_final}

\lean{Copying/CombinedLocality}{98}{104}{Combined.left_return_final}

\Needspace{12\baselineskip}
\subsection{An unbounded execution instance and exact representation}

\termdef{combined-allowance-execution}
The unbounded state type is the dependent union over all natural allowances
of closed-interface states. A transition lifts a fixed-allowance transition;
it does not add padding steps. Finality is inherited from that component,
and normality follows by inversion of a lifted step. This defines a concrete
machine and a concrete execution instance. Initialization retains the supplied
closed graph literally and places search at its unique external term input.

Here ``unbounded'' means there is no global upper bound on permitted finite
allowances. It does not mean that different padded representations have been
identified. Unlike the completed linear development, this subsystem is not
yet quotiented by padding; forward exact-cost transport is now proved in
Section~\ref{sec:copying-arithmetic-padding}, while reflection of runs from
the larger allowance remains unfinished.

\lean{Copying/CombinedExecution}{9}{9}{Combined.UnboundedState}

\lean{Copying/CombinedExecution}{10}{11}{Combined.UnboundedStep}

\lean{Copying/CombinedExecution}{13}{16}{Combined.unbounded_final_normal}

\lean{Copying/CombinedExecution}{18}{21}{Combined.unboundedMachine}

\lean{Copying/CombinedExecution}{24}{26}{Combined.execution}

Each transition preserves its allowance by construction. Transitivity
of equality proves the same for a run. Conversely, every fixed-allowance run
lifts one transition at a time. Reflection is proved by induction on an
unbounded run: inversion of its first lifted transition identifies the source
allowance and state, and the induction hypothesis reflects the remaining run.
The dependent endpoint equality records the resulting bounded state. For
endpoints already presented at the same allowance, eliminating that equality
gives an exact equivalence of the two finite-run judgements at length $n$.
Thus the wrapper introduces neither extra runs nor extra costs.

\lean{Copying/CombinedExecution}{29}{30}{Combined.step_depth}

\lean{Copying/CombinedExecution}{32}{35}{Combined.run_depth}

\lean{Copying/CombinedExecution}{38}{42}{Combined.run_unbounded}

\lean{Copying/CombinedExecution}{45}{58}{Combined.run_bounded}

\lean{Copying/CombinedExecution}{61}{68}{Combined.run_unbounded_iff}

The representation relation on unbounded states lifts a state
isomorphism at the same allowance. It is symmetric, and related states have
equivalent finality. Inverting this relation and an unbounded step reduces
one-step transport to the fixed-allowance theorem. Induction transports
unbounded runs with their exact lengths and related endpoints.

\lean{Copying/CombinedExecution}{71}{72}{Combined.Represents}

\lean{Copying/CombinedExecution}{74}{75}{Combined.represents_symm}

\lean{Copying/CombinedExecution}{77}{79}{Combined.represents_final}

\lean{Copying/CombinedExecution}{81}{88}{Combined.unbounded_step_transport}

\lean{Copying/CombinedExecution}{90}{98}{Combined.unbounded_run_transport}

For any cost preorder $Q$, a successful source run of length $n$
transports to a successful target run of the same length. Reflexivity
$Q(n,n)$ supplies the quantitative comparison. Applying this reasoning in
both directions proves machine equivalence for every cost preorder,
including exact equality of costs. Initial states of isomorphic closed
graphs have equal search modes and the same external input, so they are
related. This discharges the execution-isomorphism interface law for this
concrete subsystem.

\lean{Copying/CombinedExecution}{101}{105}{Combined.represents_refines}

\lean{Copying/CombinedExecution}{107}{108}{Combined.represents_equiv}

\lean{Copying/CombinedExecution}{111}{113}{Combined.execution_respectsIso}

The existing graph-context congruence theorem now specializes to the
concrete execution. Its proof uses the previously established composition
of observation contexts and equality up to graph isomorphism; the new
execution-isomorphism theorem supplies the formerly abstract premise.
The result preserves any already proved contextual equivalence under an
admissible graph context, with the universal cost relation. It does not
assert that any new program pair is contextually equivalent, or discharge
BPerm's operational robustness. It also does not remove the independent
padding premise from the observation-reflection theorem.

\lean{Copying/CombinedExecution}{116}{119}{Combined.equiv_context}

\Needspace{12\baselineskip}
\subsection{Every existing local execution embeds}

Each old machine step is mapped to the corresponding combined step.
For allocation, ordinary search and the raw rule map directly. Its additional
suffix-allocation constructor is handled by left serial lifting, using the
proved equality for the external input. Induction on a complete allocation
run preserves its length and literal endpoint graphs.

\lean{Copying/CombinedEmbedding}{9}{14}{Combined.allocation_step}

\lean{Copying/CombinedEmbedding}{16}{21}{Combined.allocation_run}

Dereference and assignment use their existing intrinsic search
constructors or their corresponding local-rule constructors. Assignment's
stability witness is retained. In each case induction on a run yields a
combined run with exactly the same endpoints and cost.

\lean{Copying/CombinedEmbedding}{23}{26}{Combined.dereference_step}

\lean{Copying/CombinedEmbedding}{28}{33}{Combined.dereference_run}

\lean{Copying/CombinedEmbedding}{35}{38}{Combined.assignment_step}

\lean{Copying/CombinedEmbedding}{40}{45}{Combined.assignment_run}

The atom-equality embedding handles both same-identity and distinct-
identity rules, retaining their contraction trees. Application retains the
original box body, argument and stability witness. Copying retains the original
residual tree, value and copyability witness. Each one-step case analysis
extends by run induction, again without changing any costs. These embeddings
are inclusions; they do not assert that the richer combined system has only
the runs available in an isolated local machine.

\lean{Copying/CombinedEmbedding}{47}{51}{Combined.equality_step}

\lean{Copying/CombinedEmbedding}{53}{58}{Combined.equality_run}

\lean{Copying/CombinedEmbedding}{60}{63}{Combined.application_step}

\lean{Copying/CombinedEmbedding}{65}{70}{Combined.application_run}

\lean{Copying/CombinedEmbedding}{72}{75}{Combined.copying_step}

\lean{Copying/CombinedEmbedding}{77}{82}{Combined.copying_run}

\Needspace{12\baselineskip}
\subsection{Term-wire typing throughout execution}

Both active coordinates of every local rule are term typed, by the
literal interface definitions. Induction on a rewrite preserves this
property: recursive isomorphisms preserve vertex labels; replacement and
expansion insertion preserve the labels of their inserted coordinates.
Case analysis then combines rewrite typing with the established intrinsic
search endpoint theorem. Thus both endpoints of every combined transition
are on term wires, including coordinates that have become internal.

\lean{Copying/CombinedTyping}{10}{12}{Combined.local_endpoints_term}

\lean{Copying/CombinedTyping}{15}{27}{Combined.rewrite_endpoints_term}

\lean{Copying/CombinedTyping}{29}{33}{Combined.step_endpoints_term}

A run preserves a term-typed token. The zero-length case uses its
initial typing hypothesis; in the nonempty case the first step supplies
typing at the intermediate state and induction handles the rest. Inversion
of a lifted transition gives the unbounded one-step theorem. Finally, reflect
any run from a closed program into its initial allowance and use the run
theorem. The initial input has type term by the closed-program interface.
This is an operational token invariant, not yet the stronger preservation
of all stable/copyable graph predicates required by the paper.

\lean{Copying/CombinedTyping}{36}{41}{Combined.run_term}

\lean{Copying/CombinedTyping}{43}{45}{Combined.unbounded_step_term}

\lean{Copying/CombinedTyping}{48}{53}{Combined.reachable_term}

\Needspace{12\baselineskip}
\subsection{Concrete execution regressions}

The reference program is the actual Figure~16 translation of
$\mathtt{ref}\ n$, with its distributors retained. Its safety proofs are the
translation's existing focus-free and hole-free theorems. Embedding the
previous allocation run and lifting the allowance gives a five-step
successful execution. The final state explicitly witnesses termination.
Any recursively isomorphic closed representation has exactly the same
successful costs, by the concrete representation-equivalence theorem.

\lean{Copying/CombinedExamples}{10}{14}{Combined.Examples.referenceProgram}

\lean{Copying/CombinedExamples}{16}{19}{Combined.Examples.reference_run}

\lean{Copying/CombinedExamples}{21}{22}{Combined.Examples.reference_terminates}

\lean{Copying/CombinedExamples}{25}{29}{Combined.Examples.reference_representation}

A second regression places that unchanged translated allocation in
the argument of an outer reference-operation corolla. The two serial graphs
record its source and its allocated inner reduct, and the coordinate selector
records the argument's input as inserted in that host. This is a concrete
assembled context test; no equality with the complete Figure~16 translation
of a nested reference expression is asserted.

\lean{Copying/CombinedExamples}{31}{31}{Combined.Examples.outerRef}

\lean{Copying/CombinedExamples}{32}{32}{Combined.Examples.nested}

\lean{Copying/CombinedExamples}{33}{33}{Combined.Examples.nestedAllocated}

\lean{Copying/CombinedExamples}{34}{35}{Combined.Examples.innerInput}

The outer corolla first searches its argument. The general right
serial run theorem performs the five-step allocation inside that argument.
The returned token is then handled by the outer corolla's backtrack rule,
which requests the outer rewrite. Concatenating these traces gives seven
transitions. The intermediate returned argument is not final because its
input is internal; the seven-step endpoint is not final because its mode
is rewrite. Thus the integration checks actual graph change in a context
and subsequent host search, without claiming the next allocation has run.

\lean{Copying/CombinedExamples}{37}{41}{Combined.Examples.enter_inner}

\lean{Copying/CombinedExamples}{43}{46}{Combined.Examples.return_to_outer}

\lean{Copying/CombinedExamples}{49}{54}{Combined.Examples.nested_run}

\lean{Copying/CombinedExamples}{57}{60}{Combined.Examples.inner_return_not_final}

\lean{Copying/CombinedExamples}{63}{65}{Combined.Examples.nested_endpoint_not_final}

The application regression uses a literal translated constant body,
including its unused-parameter distributor, with a stable unit argument.
The local rule's safety theorem supplies the closed-program certificates.
Embedding its seven-step opening-and-return run gives an unbounded run
and successful termination. This is the concrete application redex with an
original translated body, not a claim that the entire translated application
has already been matched through all surrounding distributors.

\lean{Copying/CombinedExamples}{68}{74}{Combined.Examples.applicationProgram}

\lean{Copying/CombinedExamples}{76}{80}{Combined.Examples.application_run}

\lean{Copying/CombinedExamples}{82}{83}{Combined.Examples.application_terminates}

At this integration milestone, the full build checked 349 jobs and the audits passed: 6 abstract,
477 graph and 830 copying declarations, with only standard Lean logical
axioms. Appendix~B.1 removes the completed non-arithmetic context integration
and its execution-isomorphism premise, while retaining the full-machine
extensions and independent padding law. No whole remaining package is declared
complete. Estimated progress by mathematical effort is about 73\%, with broad
uncertainty of 58--83\% (26 October 2026). The change reflects operational
integration and concrete execution laws, not the number of new declarations;
template robustness remains a substantial source of uncertainty.

